% ==========================================================
\chapter{Indicateurs de suivi du projet}
\label{annexe:kpi-projet}
% ==========================================================

Cette annexe présente les indicateurs clés de performance, ou KPI, retenus pour le suivi du projet \emph{5 Secondes Chrono}. 
Ces indicateurs permettent d'évaluer l'avancement du projet, la qualité technique de la solution, ainsi que la valeur d'usage attendue pour les apprenants et pour Le Carré Pro.

Les KPI sont regroupés en quatre catégories :
\begin{itemize}\setlength{\itemsep}{2pt}
    \item les KPI de pilotage projet ;
    \item les KPI de valeur et d'usage ;
    \item les KPI de performance technique ;
    \item les KPI de performance sous charge contrôlée.
\end{itemize}

Dans le cadre du projet académique, tous les indicateurs n'ont pas pu être mesurés avec le même niveau de précision. 
Les KPI de pilotage projet ont été les plus exploitables, car ils reposent sur les jalons, les exigences, les tickets Jira et les livrables. 
Les KPI techniques ont pu être observés de manière qualitative lors des tests internes et des démonstrations. 
En revanche, les KPI de valeur, d'usage et de charge nécessitent une utilisation réelle de l'application par des apprenants sur une période suffisante.

Ainsi, certains indicateurs sont considérés comme des indicateurs de suivi cible. 
Ils constituent une base pour une future phase de recette, de mise en production ou de reprise du projet.

% ==========================================================
\section{KPI de pilotage projet}
% ==========================================================

Les KPI de pilotage projet permettent de suivre la maîtrise du planning, la couverture du périmètre et la capacité de l'équipe à livrer une version démontrable du produit.

\begin{longtable}{|p{4cm}|p{3cm}|p{4cm}|p{3cm}|}
\hline
\rowcolor{etnaBlue!15}
\textbf{Indicateur} & \textbf{Objectif cible} & \textbf{Rôle dans le projet} & \textbf{Résultat observé} \\
\hline
\endfirsthead

\hline
\rowcolor{etnaBlue!15}
\textbf{Indicateur} & \textbf{Objectif cible} & \textbf{Rôle dans le projet} & \textbf{Résultat observé} \\
\hline
\endhead

Respect des jalons 
& $\geq 90\%$ 
& Vérifier la maîtrise du calendrier et des échéances principales. 
& Globalement respecté, avec des glissements et un niveau de finition variable selon les jalons. \\
\hline

Bugs critiques 
& 0 
& Garantir une version démontrable sans blocage majeur. 
& Aucun bug bloquant majeur identifié pour la démonstration finale, mais des tests complémentaires restent nécessaires. \\
\hline

Couverture fonctionnelle du MVP 
& 100\% des fonctions critiques 
& Vérifier que le cœur du produit est couvert : quiz, résultat, progression, administration minimale. 
& MVP globalement couvert sur les fonctions essentielles, avec certaines fonctions avancées reportées. \\
\hline

Dette technique maîtrisée 
& $\leq 15\%$ 
& Limiter les développements provisoires ou difficiles à maintenir. 
& Dette technique présente, notamment sur la sécurisation finale, les tests automatisés et l'industrialisation. \\
\hline

Documentation disponible 
& Documentation projet et technique livrée 
& Faciliter la compréhension, la reprise et la maintenance du projet. 
& Documentation projet disponible ; documentation technique à compléter pour une reprise complète. \\
\hline

\caption{KPI de pilotage projet}
\label{tab:kpi-pilotage-projet}
\end{longtable}

% ==========================================================
\section{KPI de valeur et d'usage}
% ==========================================================

Les KPI de valeur et d'usage permettent d'évaluer l'intérêt pédagogique de l'application et son appropriation par les utilisateurs. 
Ils sont particulièrement importants pour vérifier si l'application répond à son objectif principal : rendre l'apprentissage juridique plus régulier, plus court et plus engageant.

\begin{longtable}{|p{4.2cm}|p{3cm}|p{4.2cm}|p{3cm}|}
\hline
\rowcolor{accentGreen!15}
\textbf{Indicateur} & \textbf{Objectif cible} & \textbf{Interprétation attendue} & \textbf{Résultat observé} \\
\hline
\endfirsthead

\hline
\rowcolor{accentGreen!15}
\textbf{Indicateur} & \textbf{Objectif cible} & \textbf{Interprétation attendue} & \textbf{Résultat observé} \\
\hline
\endhead

Taux d'engagement hebdomadaire 
& $\geq 60\%$ 
& Mesurer la capacité de l'application à faire revenir les utilisateurs régulièrement. 
& Non mesurable sans mise à disposition réelle auprès d'apprenants. \\
\hline

Nombre de quiz réalisés par utilisateur 
& 3 à 5 par semaine 
& Vérifier que le format court favorise une pratique régulière. 
& Non mesurable dans le cadre du projet académique ; indicateur à suivre lors d'une future recette. \\
\hline

Taux de réussite moyen 
& 60 à 80\% 
& Évaluer si les contenus sont adaptés au niveau des apprenants. 
& Non mesurable sur un panel réel ; calcul prévu à partir des résultats de quiz. \\
\hline

Progression des apprenants 
& $\geq 70\%$ 
& Observer l'amélioration des résultats au fil des sessions. 
& Non mesurable sans historique d'usage suffisant. \\
\hline

Consultation des vidéos associées 
& À définir 
& Mesurer l'utilisation des contenus pédagogiques après le quiz. 
& Indicateur prévu pour évaluer l'intérêt des vidéos après les réponses. \\
\hline

Satisfaction utilisateur 
& $\geq 4/5$ 
& Évaluer la perception générale de l'expérience apprenant. 
& Non mesurable sans phase de test utilisateur complète. \\
\hline

\caption{KPI de valeur et d'usage}
\label{tab:kpi-valeur-usage}
\end{longtable}

\begin{tcolorbox}[
  colback=green!5,
  colframe=green!50!black,
  title={\textbf{Lecture pédagogique des KPI}},
  boxrule=0.7pt,
  enhanced,
  sharp corners,
  breakable
]
Les KPI de valeur ne doivent pas être lus isolément. 
Par exemple, un taux de réussite élevé peut indiquer une bonne progression, mais aussi des questions trop simples si l'engagement reste faible. 
À l'inverse, un taux de réussite moyen associé à un engagement fort peut traduire une dynamique d'apprentissage positive.
\end{tcolorbox}

% ==========================================================
\section{KPI de performance technique}
% ==========================================================

Les KPI de performance technique permettent de vérifier que l'application est suffisamment fluide, stable et exploitable pour une première version.

\begin{longtable}{|p{5cm}|p{3cm}|p{4cm}|p{2.5cm}|}
\hline
\rowcolor{orange!20}
\textbf{Indicateur} & \textbf{Objectif cible} & \textbf{Rôle} & \textbf{Résultat observé} \\
\hline
\endfirsthead

\hline
\rowcolor{orange!20}
\textbf{Indicateur} & \textbf{Objectif cible} & \textbf{Rôle} & \textbf{Résultat observé} \\
\hline
\endhead

Temps de chargement initial de l'application 
& $< 3$ secondes 
& Garantir une entrée rapide dans l'application. 
& Observé qualitativement comme acceptable en démonstration ; mesure automatisée à prévoir. \\
\hline

Temps de réponse des écrans critiques 
& $< 2$ secondes 
& Préserver la fluidité du quiz, des résultats et de la progression. 
& Parcours démontrable sans lenteur bloquante ; tests de performance à compléter. \\
\hline

Temps de réponse moyen de l'API 
& $< 500$ ms 
& Vérifier la rapidité des échanges entre frontend et backend. 
& Non mesuré de manière systématique ; indicateur à instrumenter avec logs ou monitoring. \\
\hline

Taux d'erreurs applicatives 
& $< 1\%$ 
& Suivre la stabilité globale de la solution. 
& Non mesuré automatiquement ; les erreurs ont été traitées lors des tests internes. \\
\hline

Disponibilité de l'application 
& $\geq 99\%$ 
& Évaluer la fiabilité du service dans un contexte d'exploitation. 
& Non applicable hors production réelle. \\
\hline
\caption{KPI de performance technique}
\label{tab:kpi-performance-technique}
\end{longtable}

% ==========================================================
\section{KPI de performance sous charge contrôlée}
% ==========================================================

Dans le cadre académique du projet, les tests de charge ne visent pas à garantir une montée en charge industrielle. 
Ils permettent surtout de vérifier que la solution reste stable dans un contexte d'usage pilote ou de démonstration.

\begin{longtable}{|p{5cm}|p{3cm}|p{4cm}|p{2.5cm}|}
\hline
\rowcolor{purple!15}
\textbf{Indicateur} & \textbf{Objectif cible} & \textbf{Rôle} & \textbf{Résultat observé} \\
\hline
\endfirsthead

\hline
\rowcolor{purple!15}
\textbf{Indicateur} & \textbf{Objectif cible} & \textbf{Rôle} & \textbf{Résultat observé} \\
\hline
\endhead

Nombre d'utilisateurs simultanés supportés 
& $\geq 50$ 
& Vérifier la tenue de l'application en usage simultané limité. 
& Non testé dans le cadre académique ; objectif à conserver pour une future phase de recette technique. \\
\hline

Dégradation des temps de réponse sous charge 
& $< +30\%$ 
& Observer l'impact de la charge sur la fluidité. 
& Non mesuré ; nécessite un environnement de test dédié. \\
\hline

Stabilité de l'application sous charge 
& Aucun crash 
& Vérifier la robustesse minimale du système. 
& Non validé sous charge réelle ; tests à prévoir avant mise en production. \\
\hline

\caption{KPI de performance sous charge contrôlée}
\label{tab:kpi-charge}
\end{longtable}

% ==========================================================
\section{Synthèse de suivi des KPI}
% ==========================================================

Le tableau suivant permet de faire une synthèse globale de l'état des KPI en fin de projet.

\begin{table}[H]
\centering
\renewcommand{\arraystretch}{1.25}
\small
\begin{tabular}{|p{4cm}|p{3cm}|p{3cm}|p{4cm}|}
\hline
\rowcolor{gray!20}
\textbf{Famille de KPI} & \textbf{Mesurable pendant le projet} & \textbf{Statut global} & \textbf{Commentaire} \\
\hline

Pilotage projet 
& Oui 
& Exploitable 
& Les jalons, exigences, tickets Jira et livrables ont permis un suivi global du projet. \\
\hline

Valeur et usage 
& Partiellement 
& À mesurer en recette 
& Ces KPI nécessitent un panel d'apprenants et une période d'utilisation réelle. \\
\hline

Performance technique 
& Partiellement 
& Observé qualitativement 
& Les parcours ont été testés en démonstration, mais sans instrumentation complète. \\
\hline

Charge contrôlée 
& Non 
& À prévoir 
& Les tests de charge nécessitent un environnement dédié et une campagne spécifique. \\
\hline

\end{tabular}
\caption{Synthèse globale du suivi des KPI}
\label{tab:synthese-kpi}
\end{table}

% ==========================================================
\section{Bilan d'utilisation des KPI}
% ==========================================================

Les KPI ont été définis comme des repères de pilotage et non comme de simples objectifs chiffrés. 
Ils permettent de relier l'avancement du projet à des critères concrets : couverture du MVP, qualité technique, stabilité, engagement utilisateur et valeur pédagogique.

Dans le cadre du projet académique, tous les KPI n'ont pas pu être mesurés avec le même niveau de précision. 
Les indicateurs de pilotage projet sont les plus exploitables, car ils reposent sur des éléments disponibles : exigences, jalons, tickets Jira, livrables et état du MVP. 
Les indicateurs techniques ont principalement été observés lors des tests internes et des démonstrations, mais ils nécessiteraient une instrumentation plus complète pour une exploitation réelle.

Les KPI de valeur et d'usage restent principalement prospectifs. 
Ils ne peuvent être interprétés correctement qu'après une phase d'utilisation par des apprenants réels. 
Ils sont néanmoins essentiels pour la suite du projet, car ils permettront de vérifier si l'application remplit réellement sa promesse : favoriser un apprentissage juridique court, régulier et engageant.

Enfin, les KPI de charge contrôlée doivent être considérés comme des critères de préparation à une future mise en production. 
Ils permettront d'évaluer la robustesse de l'application dans un contexte d'usage simultané, mais n'ont pas vocation à être validés uniquement dans le cadre d'une démonstration académique.

\begin{tcolorbox}[
  colback=blue!5,
  colframe=blue!70!black,
  title={\textbf{Lecture globale des KPI}},
  boxrule=0.7pt,
  enhanced,
  sharp corners,
  breakable
]
Les KPI définis dans cette annexe constituent une base de pilotage pour la suite du projet. 
Ils montrent ce qui a pu être suivi pendant le cadre académique, mais aussi ce qui devra être mesuré lors d'une future recette utilisateur, d'une phase pilote ou d'une mise en production.
\end{tcolorbox}